(1) 取 \(\mathbf H=H_z(x,t)\hat{\mathbf z}\)、\(\mathbf E=E_y(x,t)\hat{\mathbf y}\)。忽略位移电流后 Ampère 定律给 \(J_y=-\partial_xH_z=\sigma E_y\)；Faraday 定律给 \(\partial_xE_y=-\mu\partial_tH_z\)。消去 \(E_y\) 得 \(\partial_x^2H_z=\mu\sigma\partial_tH_z\)，这是扩散而非无耗波动方程。取 \(\widetilde H_z=H_0e^{-kx}\)，由 \(k^2=-i\mu\sigma\omega\) 且要求 \(\Re k>0\)，选 \(k=(1-i)/\delta\)，\(\delta=\sqrt{2/(\mu\sigma\omega)}\)。实际场为 \(H_z=H_0e^{-x/\delta}\cos(\omega t-x/\delta)\)，在深度 \(\delta\) 振幅降为表面的 \(e^{-1}\)，同时相位滞后 \(1\) 弧度。

(2) 由 \(\nabla\times\mathbf H=\mathbf J\) 得 \(\widetilde J_y=kH_0e^{-kx}\)，由局域 Ohm 定律得 \(\widetilde E_y=(k/\sigma)H_0e^{-kx}\)。因此 \(Z_s=k/\sigma=(1-i)/(\sigma\delta)\)，实部正而虚部在 \(e^{-i\omega t}\) 约定下为负。表面实际电场为 \(E_y(0,t)=\sqrt2 H_0/(\sigma\delta)\cos(\omega t+\pi/4)\)，所以电场时间相位比表面磁场超前 \(\pi/4\)。再独立代回 Faraday 方程：\(\partial_x\widetilde E_y=-k^2\widetilde H_z/\sigma=i\mu\omega\widetilde H_z=-\partial_t\widetilde B_z\) 的相量形式，符号与取向完全一致。

(3) 在表面 \(\mathbf E\parallel\hat{\mathbf y}\)、\(\mathbf H\parallel\hat{\mathbf z}\)，故 \(\mathbf E\times\mathbf H\) 指向导体内部 \(+\hat{\mathbf x}\)。平均输入功率通量 \(\langle S_x(0)\rangle=\tfrac12\Re[\widetilde E_y(0)\widetilde H_z(0)^*]=\tfrac12|H_0|^2\Re Z_s=|H_0|^2/(2\sigma\delta)\)。另一方面 \(\langle p_J(x)\rangle=|\widetilde J_y(x)|^2/(2\sigma)=|k|^2|H_0|^2e^{-2x/\delta}/(2\sigma)\)。积分 \(\int_0^\infty e^{-2x/\delta}dx=\delta/2\)，利用 \(|k|^2=2/\delta^2\) 得 \(\int_0^\infty\langle p_J\rangle dx=|H_0|^2/(2\sigma\delta)\)。两种完全不同的能量账本相等，表明没有漏掉表面功率或虚假的能量源。

(4) 磁能密度的周期平均为 \(\mu|\widetilde H_z|^2/4\)，故每单位面积储存的平均磁能 \(W_m/A=\mu|H_0|^2\int_0^\infty e^{-2x/\delta}dx/4=\mu|H_0|^2\delta/8\)。因为 \(\mu\sigma\omega\delta^2=2\)，有 \(2\omega W_m/A=|H_0|^2/(2\sigma\delta)=|\Im Z_s|\,|H_0|^2/2\)，对应表面感性无功功率的大小；这里虚部符号取决于相量时间约定。趋肤深度随 \(\omega^{-1/2}\) 缩小，表面电阻随 \(\omega^{1/2}\) 增大。必须同时满足 \(\sigma\gg\omega\varepsilon\)、样品厚度远大于 \(\delta\) 且局域 Ohm 定律适用；对有限厚样品不能直接令 \(\omega\to0\) 后仍使用半空间公式。
